% ============================================================
%  CHƯƠNG 4: BÀI ĐĂNG & TÌM KIẾM
%  Phụ trách: [Thành viên 3]
%  Packages: ui/main/home, ui/main/search, ui/postdetail,
%            ui/main/myposts
%  Models:   data/model/post/*, data/model/search/*,
%            data/model/rating/*, data/model/location/*
%  Ghi chú: File này được \input{} từ main.tex
% ============================================================

\chapter{Bài đăng và Tìm kiếm}

Chương này trình bày phân hệ trung tâm của ứng dụng Trọ Tốt: hiển thị, tìm kiếm,
đăng tải và đánh giá bài đăng phòng trọ. Nhóm viết toàn bộ phân hệ bằng
\tech{Java}, bám sát kiến trúc MVVM ở Chương 1. Giao diện dùng các thành phần
chuẩn của Android (\code{TextView}, \code{RecyclerView} -- thành phần kế thừa và
thay thế ListView/GridView). Tầng dữ liệu gọi API qua \tech{Retrofit}.

% ============================================================
\section{Màn hình Trang chủ (Home Feed)}
% ============================================================

\subsection{Yêu cầu chức năng}
Trang chủ là màn hình đầu tiên người dùng nhìn thấy sau khi mở ứng dụng, với hai
khối nội dung chính:
\begin{itemize}
    \item \textbf{Tin đăng mới nhất:} hiển thị 4 bài đăng đã duyệt gần nhất,
    giúp người dùng nắm bắt nhanh nguồn cung phòng trọ mới.
    \item \textbf{Gợi ý dành cho bạn (Recommendation):} hệ thống gợi ý bài
    đăng theo lịch sử tương tác của từng người dùng, có phân trang để tải thêm.
\end{itemize}
Người dùng chưa đăng nhập vẫn xem được tin mới nhất; khối gợi ý cá nhân hóa
yêu cầu đã đăng nhập (API gắn JWT).

\subsection{Thiết kế giao diện}
File \code{fragment\_home.xml} định nghĩa giao diện với các thành phần:
\begin{itemize}
    \item \code{TextView}: tiêu đề ứng dụng trên thanh công cụ và tiêu đề từng
    khối nội dung (``Tin đăng mới nhất'', ``Gợi ý dành cho bạn'').
    \item Hai \code{RecyclerView} (\code{rvLatestPosts}, \code{rvRecommendations})
    dùng chung \code{PostAdapter} hiển thị thẻ bài đăng
    (\code{item\_post\_card.xml}): ảnh đại diện tải bằng \tech{Glide}, tiêu đề,
    giá thuê, diện tích và địa chỉ.
    \item \code{ProgressBar} hiển thị trạng thái đang tải;
    \code{MaterialButton} ``Thử lại'' khi lỗi mạng và ``Tải thêm''
    (\code{btnLoadMore}) cho phân trang gợi ý.
    \item Toàn bộ nội dung đặt trong \code{NestedScrollView} để cuộn mượt cả
    hai danh sách.
\end{itemize}

% TODO: Bỏ comment sau khi chụp màn hình và lưu ảnh vào images/outputs/
% \begin{figure}[H]
%     \centering
%     \includegraphics[width=0.45\textwidth]{images/outputs/c3_home_feed.png}
%     \caption{Màn hình Trang chủ với tin mới nhất và gợi ý cá nhân hóa}
%     \label{fig:c3_home_feed}
% \end{figure}

\subsection{Luồng xử lý}
\begin{enumerate}
    \item \code{HomeFragment} khởi chạy, gọi \code{loadLatestPosts()} và
    \code{loadRecommendations()} trên \code{HomeViewModel}.
    \item ViewModel phát trạng thái \code{Resource.loading()} -- UI hiện
    \code{ProgressBar}.
    \item \code{PostRepository} thực thi request Retrofit trên luồng IO
    (\code{subscribeOn(Schedulers.io())}).
    \item Repository bọc kết quả vào \code{Resource<List<Post>>} rồi đẩy vào
    \code{LiveData} trên Main Thread.
    \item Fragment (đang observe) cập nhật adapter; nếu lỗi thì hiển thị thông
    báo cùng nút ``Thử lại''.
    \item Khi người dùng nhấn ``Tải thêm'', ViewModel gọi
    \code{loadMoreRecommendations()}: tăng \code{page}, gửi kèm
    \code{recommendationLogId} của phiên gợi ý để backend ghi nhận chuỗi
    tương tác, sau đó nối (append) kết quả mới vào danh sách hiện có.
\end{enumerate}

\subsection{Triển khai trên Android}
\begin{itemize}
    \item \textbf{Lớp chính:} \code{HomeFragment}, \code{HomeViewModel}
    (\code{@HiltViewModel}), \code{PostAdapter} thuộc package
    \code{ui/main/home}; \code{PostRepository} thuộc \code{data/repository}.
    \item \textbf{API sử dụng:}
    \begin{itemize}
        \item \code{GET api/posts/latest?limit=4} $\rightarrow$
        \code{ResponseData<List<Post>>}
        \item \code{GET api/recommend?page\&pageSize\&logId} $\rightarrow$
        \code{RecommendationResponse} (kèm \code{pagination.hasMore} và
        \code{recommendationLogId})
    \end{itemize}
    \item ViewModel giữ toàn bộ \textbf{trạng thái phân trang} (trang hiện tại,
    danh sách tích lũy \code{allRecommendations}, cờ \code{isLoadingMore},
    \code{hasMoreRecommendations}). Nhờ đó dữ liệu không mất khi xoay màn hình.
\end{itemize}

\textit{Đáp ứng CLO1 (chức năng + giao diện hoàn chỉnh bằng Java) và CLO2: sử
dụng TextView, RecyclerView và gọi API qua Retrofit.}

% ============================================================
\section{Tìm kiếm bài đăng}
% ============================================================

\subsection{Yêu cầu chức năng}
Người dùng nhập từ khóa tự nhiên (ví dụ ``phòng trọ gần đại học có gác'') và
thu hẹp kết quả bằng bộ lọc: tỉnh/thành phố, quận/huyện, phường/xã, khoảng
giá, khoảng diện tích, tình trạng nội thất. Kết quả trả về theo trang
(20 bài/trang), tải tiếp khi cuộn hết. Hệ thống đồng thời ghi nhận tương tác
tìm kiếm (lượt nhấn kết quả, phản hồi chất lượng) phục vụ cải thiện mô hình
gợi ý.

\subsection{Thiết kế giao diện}
Giao diện \code{fragment\_search.xml} gồm:
\begin{itemize}
    \item \code{TextInputEditText} nhập từ khóa kèm nút tìm kiếm.
    \item \code{Chip} (Material) hiển thị bộ lọc đang áp dụng;
    \code{MaterialButton} mở hộp thoại bộ lọc.
    \item \code{RecyclerView} hiển thị danh sách kết quả, tái sử dụng
    \code{PostAdapter} của Trang chủ.
    \item \code{TextView} thông báo khi không có kết quả hoặc khi lỗi.
\end{itemize}
\code{SearchFilterBottomDialog} (kế thừa \code{BottomSheetDialogFragment})
đảm nhận bộ lọc: ô nhập khoảng giá/diện tích, danh sách chọn địa giới hành
chính và \code{ChipGroup} chọn tình trạng nội thất (\code{chipGroupInterior}).

% TODO: Bỏ comment sau khi chụp màn hình tìm kiếm + bottom sheet bộ lọc
% \begin{figure}[H]
%     \centering
%     \begin{subfigure}{0.45\textwidth}
%         \includegraphics[width=\textwidth]{images/outputs/c3_search_results.png}
%         \caption{Kết quả tìm kiếm}
%     \end{subfigure}
%     \hfill
%     \begin{subfigure}{0.45\textwidth}
%         \includegraphics[width=\textwidth]{images/outputs/c3_search_filter.png}
%         \caption{Bộ lọc tìm kiếm}
%     \end{subfigure}
%     \caption{Màn hình tìm kiếm bài đăng}
%     \label{fig:c3_search}
% \end{figure}

\subsection{Luồng xử lý}
\begin{enumerate}
    \item Người dùng nhập từ khóa; \code{SearchViewModel} kiểm tra hợp lệ (từ
    khóa không rỗng) trước khi tìm.
    \item ViewModel gom từ khóa và toàn bộ bộ lọc đang chọn vào
    \code{SearchParams} qua hàm \code{buildSearchParams()}, đồng thời reset
    phân trang về trang 1.
    \item \code{PostRepository.search()} gọi API tìm kiếm. Backend xử lý bằng
    hybrid vector search -- phía Android chỉ gọi API và hiển thị kết quả.
    \item ViewModel đẩy kết quả về UI qua
    \code{LiveData<Resource<List<Post>>>} và giữ lại \code{searchLogId} trong
    response.
    \item Người dùng nhấn một kết quả: ứng dụng gọi
    \code{logClick(searchLogItemId)} ghi nhận lượt nhấn gắn với
    \code{searchLogId}, rồi điều hướng sang màn hình chi tiết.
    \item Cuộn đến cuối danh sách và còn trang kế (\code{hasNextPage}),
    ViewModel gọi \code{loadNextPage()} nối thêm kết quả.
\end{enumerate}

\subsection{Triển khai trên Android}
\begin{itemize}
    \item \textbf{Lớp chính:} \code{SearchFragment}, \code{SearchViewModel},
    \code{SearchFilterBottomDialog} (package \code{ui/main/search}); model
    \code{SearchParams}, \code{SearchResponse},
    \code{SearchInteractionRequest} (package \code{data/model/search}).
    \item \textbf{API sử dụng:}
    \begin{itemize}
        \item \code{GET api/search} với các query:
        \code{query}, \code{city}, \code{district}, \code{ward},
        \code{priceMin}, \code{priceMax}, \code{acreageMin}, \code{acreageMax},
        \code{interiorCondition}, \code{page}, \code{pageSize}.
        \item \code{POST api/search/click} -- ghi nhận lượt nhấn kết quả.
        \item \code{POST api/search/feedback} -- gửi phản hồi chất lượng
        tìm kiếm (hữu ích hay không, vấn đề gặp phải, bình luận).
    \end{itemize}
    \item Bottom sheet bộ lọc dùng chung ViewModel với Fragment cha: khởi tạo
    \code{ViewModelProvider} với \code{getParentFragment()}. Trạng thái lọc
    vì vậy không mất khi đóng/mở bottom sheet.
\end{itemize}

\textit{Đáp ứng CLO2: thành phần nhập liệu (TextInputEditText, ChipGroup),
RecyclerView hiển thị danh sách và gọi API tìm kiếm có tham số.}

% ============================================================
\section{Chi tiết bài đăng}
% ============================================================

\subsection{Yêu cầu chức năng}
Màn hình chi tiết hiển thị đầy đủ thông tin một phòng trọ: thư viện ảnh/video,
tiêu đề, giá thuê, mô tả, bảng thông số (diện tích, nội thất, địa chỉ), bản đồ
vị trí, thông tin chủ trọ và khối đánh giá. Người dùng lưu tin, xem số điện
thoại chủ trọ (yêu cầu đăng nhập) và gửi đánh giá ngay trên màn hình này.

\subsection{Thiết kế giao diện}
Giao diện \code{fragment\_post\_detail.xml} gồm:
\begin{itemize}
    \item \code{ViewPager2} (\code{vpImages}) kết hợp thư viện
    \tech{DotsIndicator} tạo gallery ảnh trượt ngang;
    \code{PostImageAdapter} tải ảnh bằng \tech{Glide}.
    \item Các \code{TextView}: tiêu đề, giá (định dạng tiền tệ
    \code{NumberFormat} locale vi-VN), mô tả rút gọn 3 dòng kèm nút
    ``Xem thêm/Thu gọn''.
    \item Bảng thông số dựng động từ các hàng
    \code{item\_post\_detail\_row.xml}, mỗi hàng là cặp \code{TextView}
    nhãn--giá trị (vai trò tương đương bảng dữ liệu dạng lưới).
    \item \code{WebView} (\code{wvMap}) nhúng Google Maps Embed hiển thị vị trí
    phòng trọ theo địa chỉ.
    \item Khối chủ trọ: ảnh đại diện (Glide), tên, khu vực, ngày tham gia,
    nút xem SĐT và nút Chat.
    \item \code{FloatingActionButton} (\code{fabSave}) lưu/bỏ lưu tin với icon
    bookmark đổi trạng thái.
    \item Khối đánh giá: \code{RatingBar} thống kê, \code{RecyclerView}
    (\code{rvRatings}) danh sách nhận xét, form gửi đánh giá gồm
    \code{RatingBar} chọn sao và \code{EditText} bình luận.
\end{itemize}

% TODO: Bỏ comment sau khi chụp màn hình chi tiết bài đăng
% \begin{figure}[H]
%     \centering
%     \includegraphics[width=0.45\textwidth]{images/outputs/c3_post_detail.png}
%     \caption{Màn hình chi tiết bài đăng}
%     \label{fig:c3_post_detail}
% \end{figure}

\subsection{Luồng xử lý}
\begin{enumerate}
    \item Fragment nhận \code{postId} qua \code{Bundle} khi điều hướng.
    \item \code{PostDetailViewModel} tải song song: chi tiết bài đăng, thống kê
    đánh giá, đánh giá của bản thân, danh sách nhận xét và trạng thái đã lưu.
    \item Dữ liệu đổ về qua các \code{LiveData} riêng biệt; mỗi khối UI tự
    cập nhật độc lập (ảnh, bảng thông số, bản đồ, đánh giá).
    \item \textbf{Bảo vệ quyền riêng tư:} ứng dụng che số điện thoại chủ trọ
    dạng \code{09xx***xxx} (\code{getMaskedPhone()}). Số đầy đủ chỉ hiện khi
    người dùng đã đăng nhập và chủ động nhấn ``Hiện số''; hệ thống ghi nhật ký
    liên hệ ngay tại thời điểm đó.
    \item Nhấn \code{fabSave}: ViewModel gọi \code{toggleSavePost(postId)},
    cập nhật icon theo kết quả và hiện Toast xác nhận.
\end{enumerate}

\subsection{Triển khai trên Android}
\begin{itemize}
    \item \textbf{Lớp chính:} \code{PostDetailFragment},
    \code{PostDetailViewModel}, \code{PostImageAdapter},
    \code{PostRatingAdapter} (package \code{ui/postdetail}).
    \item \textbf{API sử dụng:}
    \begin{itemize}
        \item \code{GET api/posts/\{postId\}} -- chi tiết bài đăng
        (\code{PostDetail} gồm \code{MultimediaFile}, \code{PostOwner}).
        \item \code{POST api/customer/saved-posts/\{postId\}} /
        \code{DELETE api/customer/saved-posts/\{postId\}} -- lưu/bỏ lưu tin;
        \code{GET .../check} -- kiểm tra trạng thái đã lưu.
        \item \code{POST api/interactions/contact} -- ghi nhật ký khi người
        dùng xem số điện thoại liên hệ (\code{ContactLogRequest}).
    \end{itemize}
    \item Các API lưu tin và ghi nhật ký đều yêu cầu xác thực;
    \code{AuthInterceptor} tự động gắn JWT vào header như mô tả ở Chương 1.
\end{itemize}

\textit{Đáp ứng CLO2: đa dạng thành phần giao diện (TextView, ViewPager2,
WebView, RecyclerView), gọi nhiều API và bảo vệ quyền riêng tư số điện
thoại bằng cơ chế che số kết hợp yêu cầu đăng nhập.}

% ============================================================
\section{Đăng và quản lý bài đăng của tôi}
% ============================================================

\subsection{Yêu cầu chức năng}
Phân hệ dành riêng cho vai trò \textbf{chủ trọ} -- minh chứng trực tiếp cho
yêu cầu phân quyền của học phần:
\begin{itemize}
    \item Đăng tin mới kèm nhiều ảnh (mỗi ảnh $\le$ 5MB) và một video
    ($\le$ 25MB); nhập đủ tiêu đề, mô tả, giá, diện tích, địa chỉ, tình trạng
    nội thất.
    \item Xem danh sách tin của mình theo bộ lọc trạng thái, cuộn vô tận.
    \item Chỉnh sửa, ẩn/hiện tin đã đăng và theo dõi tiến độ kiểm duyệt.
\end{itemize}

\subsection{Thiết kế giao diện}
\begin{itemize}
    \item \code{fragment\_myposts.xml}: dùng \code{RecyclerView}
    (\code{rvMyPosts}) với \code{LinearLayoutManager}, bọc trong
    \code{SwipeRefreshLayout} hỗ trợ kéo làm mới; thanh công cụ có nút tạo
    tin mới. Mỗi thẻ tin hiển thị trạng thái duyệt bằng nhãn màu
    (\code{TextView}).
    \item \code{fragment\_post\_create.xml}: form nhập liệu với bộ đếm ký tự,
    danh sách ảnh đã chọn (ảnh đầu mang nhãn ``ảnh bìa''), thẻ xem trước
    video, và hộp thoại chọn địa chỉ hợp nhất (tỉnh $\rightarrow$ quận
    $\rightarrow$ phường) dùng \code{AutoCompleteTextView}.
\end{itemize}

% TODO: Bỏ comment sau khi chụp màn hình quản lý tin + form đăng tin
% \begin{figure}[H]
%     \centering
%     \begin{subfigure}{0.45\textwidth}
%         \includegraphics[width=\textwidth]{images/outputs/c3_my_posts.png}
%         \caption{Quản lý tin của tôi}
%     \end{subfigure}
%     \hfill
%     \begin{subfigure}{0.45\textwidth}
%         \includegraphics[width=\textwidth]{images/outputs/c3_post_create.png}
%         \caption{Đăng tin mới}
%     \end{subfigure}
%     \caption{Phân hệ đăng và quản lý bài đăng}
%     \label{fig:c3_my_posts}
% \end{figure}

\subsection{Luồng xử lý}
\textbf{Đăng tin mới:}
\begin{enumerate}
    \item Người dùng chọn ảnh/video qua \code{ActivityResultLauncher}
    (\code{GetContent}); ứng dụng kiểm tra kích thước tệp ngay tại client
    trước khi chấp nhận.
    \item Chọn địa giới hành chính: danh sách tỉnh/quận đọc từ
    \code{provinces.json} trong \code{assets}; danh sách phường lấy động qua
    API \code{GET api/location/wards/\{districtId\}}.
    \item Khi nhấn Đăng, \code{PostFormViewModel} đóng gói toàn bộ trường
    thành \code{multipart/form-data} (mỗi trường text là một
    \code{RequestBody}, ảnh/video là \code{MultipartBody.Part}) rồi gọi API.
    \item Tin mới mang trạng thái chờ duyệt; màn hình ``Tin của tôi'' hiển thị
    kết quả kiểm duyệt.
\end{enumerate}
\textbf{Quản lý tin:} danh sách phân trang theo con trỏ (cursor-based). Khi
\code{OnScrollListener} phát hiện đã cuộn tới cuối và \code{hasMore()} còn
\code{true}, ViewModel gọi trang kế với \code{cursor} là ID bài cuối cùng.
Cách này tránh trùng lặp dữ liệu khi có tin mới chen vào giữa hai lần tải.

\subsection{Triển khai trên Android}
\begin{itemize}
    \item \textbf{Lớp chính:} màn hình quản lý gồm \code{MyPostsFragment},
    \code{MyPostsViewModel} và \code{MyPostsAdapter}; form đăng/sửa tin gồm
    \code{PostCreateFragment}, \code{PostEditFragment} và
    \code{PostFormViewModel} (package \code{ui/main/myposts}).
    \item \textbf{API sử dụng:}
    \begin{itemize}
        \item \code{GET api/posts/me?status\&cursor\&limit} -- danh sách tin
        của tôi (\code{MyPostsResponse}).
        \item \code{POST api/posts} (multipart) -- tạo tin;
        \code{PUT api/posts/\{postId\}} (multipart) -- cập nhật tin, kèm
        trường \code{oldFiles} liệt kê tệp giữ lại.
        \item \code{POST api/posts/\{postId\}/hide} / \code{.../unhide} --
        ẩn/hiện tin.
        \item \code{GET api/posts/me/\{postId\}} -- lấy chi tiết tin của mình
        để đổ vào form chỉnh sửa.
    \end{itemize}
    \item Trạng thái bài đăng (\code{PostStatus}: chờ duyệt, đã duyệt, bị từ
    chối, đã ẩn) là enum dùng chung với phân hệ kiểm duyệt trình bày ở
    Chương 6.
\end{itemize}

\textit{Đáp ứng CLO1 và CLO2: chức năng nghiệp vụ hoàn chỉnh cho vai trò chủ
trọ (phân quyền), upload đa phương tiện qua API multipart, RecyclerView với
phân trang con trỏ.}

% ============================================================
\section{Đánh giá chủ trọ}
% ============================================================

\subsection{Yêu cầu chức năng}
Người thuê đã đăng nhập đánh giá bài đăng theo thang 5 sao kèm bình luận.
Mỗi người chỉ có một đánh giá trên mỗi tin và chỉ xóa được đánh giá của mình.
Mọi người dùng xem được điểm trung bình, tổng lượt đánh giá và danh sách
nhận xét.

\subsection{Thiết kế giao diện}
Khối đánh giá nằm trong màn hình chi tiết bài đăng:
\begin{itemize}
    \item Thống kê: \code{TextView} điểm trung bình (1 chữ số thập phân),
    \code{RatingBar} hiển thị sao và \code{TextView} số lượt đánh giá.
    \item Danh sách nhận xét: \code{RecyclerView} (\code{rvRatings}) với
    \code{item\_rating.xml} gồm tên người đánh giá, số sao và bình luận.
    \item Form gửi: \code{RatingBar} chọn sao (mặc định 5),
    \code{EditText} nhập bình luận, \code{MaterialButton} gửi. Với người đã
    đánh giá, UI thay form bằng khối ``đánh giá của tôi'' kèm nút xóa (có
    \code{AlertDialog} xác nhận). Người chưa đăng nhập thấy lời nhắc đăng nhập.
\end{itemize}

\subsection{Luồng xử lý}
\begin{enumerate}
    \item Khi mở chi tiết bài, ViewModel tải đồng thời thống kê
    (\code{RatingStats}), đánh giá của tôi và trang nhận xét đầu tiên.
    \item UI phân nhánh theo ba trạng thái: chưa đăng nhập / đã đánh giá /
    chưa đánh giá -- mỗi thời điểm chỉ hiện một khối tương ứng.
    \item Gửi đánh giá: client kiểm tra số sao $\ge$ 1, gọi API, sau đó tải
    lại thống kê và danh sách để đồng bộ.
    \item Danh sách nhận xét phân trang theo con trỏ thời gian
    (\code{cursor} là mốc ngày của nhận xét cuối).
\end{enumerate}

\subsection{Triển khai trên Android}
\begin{itemize}
    \item \textbf{Lớp chính:} \code{PostRatingAdapter},
    \code{PostDetailViewModel}; model \code{Rating}, \code{RatingStats},
    \code{AddRatingRequest}, \code{RatingListResponse}
    (package \code{data/model/rating}).
    \item \textbf{API sử dụng:}
    \begin{itemize}
        \item \code{GET api/customer/posts/\{postId\}/rate-avg} -- thống kê
        điểm trung bình và số lượt.
        \item \code{GET api/customer/posts/\{postId\}/rates?limit\&cursor} --
        danh sách nhận xét.
        \item \code{POST api/customer/posts/\{postId\}/rate} -- gửi đánh giá
        (yêu cầu đăng nhập).
        \item \code{GET} / \code{DELETE api/customer/posts/\{postId\}/rate} --
        xem / xóa đánh giá của tôi.
    \end{itemize}
\end{itemize}

\textit{Đáp ứng CLO2: thành phần RatingBar, RecyclerView, EditText; gọi API
CRUD có xác thực; chỉ chủ sở hữu đánh giá mới xóa được đánh giá của mình
(bảo vệ dữ liệu cá nhân).}

% ============================================================
\section{Kết chương}
% ============================================================
Phân hệ Bài đăng \& Tìm kiếm khép kín vòng đời một tin đăng: chủ trọ đăng tải,
người thuê khám phá qua trang chủ/tìm kiếm, xem chi tiết và phản hồi bằng đánh
giá. Phân hệ dùng xuyên suốt các thành phần giao diện chuẩn Android, giao tiếp
với backend hoàn toàn qua API REST có xác thực JWT, và phân tách rõ vai trò
người thuê -- chủ trọ đúng yêu cầu phân quyền của học phần.
